% =====================================================================
%  Extra packages + custom macros for the ISFA mémoire
%  (Quarto already loads: hyperref, xcolor, amsmath, graphicx, geometry,
%   fontenc/inputenc, natbib via cite-method. Do NOT reload those.)
% =====================================================================

% --- Languages: French main + English for the abstract / summary -----
\usepackage[english, french]{babel}

% --- Maths & symbols --------------------------------------------------
\usepackage{amsfonts, amssymb}

% --- Title-page graphics ---------------------------------------------
\usepackage{tikz}

% --- Tables & layout helpers -----------------------------------------
\usepackage{multirow, multicol, array}
\usepackage{wrapfig}
\usepackage{subfig}
\usepackage{eurosym}
\usepackage{todonotes}

% Don't hyphenate too aggressively (default 50, max 10000)
\hyphenpenalty=400
\setlength{\columnsep}{.35in}

% Footnotes restart per page, printed as symbols (as in the original)
\usepackage[stable]{footmisc}

% =====================================================================
%  Custom macros (from the original main-file.tex).
%  \providecommand avoids clashes with Quarto's own definitions.
% =====================================================================
\providecommand{\cyan}[1]{{\color{cyan}#1}}
\providecommand{\blue}[1]{{\color{blue}#1}}
\providecommand{\red}[1]{\textbf{{\color{red}#1}}}
\providecommand{\pkg}{\textbf}
\providecommand{\sigle}{\textsc}
\providecommand{\code}{\texttt}
\providecommand{\soft}{\textsf}
\providecommand{\myskip}{\vspace{\parskip}}
\providecommand{\expo}{\textsuperscript}
\providecommand{\HRule}{\noindent\rule{\textwidth}{1pt}}
\providecommand{\II}{\mbox{\large 1\hskip -0.353em 1}}
\providecommand{\abs}[1]{\lvert#1\rvert}
\providecommand{\norm}[1]{\lVert#1\rVert}
\providecommand{\R}{\mathbb{R}}
\providecommand{\card}{\textrm{Card}}
\providecommand{\logit}{\textrm{logit}}
\providecommand{\probit}{\textrm{probit}}
\providecommand{\boldit}[1]{\textbf{\textit{#1}}}
\providecommand{\fillin}[1]{\makebox[#1]{\dotfill}}

% A bit of breathing room between paragraphs, like the original
\setlength{\parskip}{0.5\baselineskip}
